
From part (1) of this exercise, we know that the moment generating function $ M_{Y_{n} } (t) $exists for all $ t \in
\R$. By Theorem 10.9 (Chernoff's estimate), for $ t \in [0, t_0]$ ($ t_0 = \infty$ for this case)
we get that 

\[
    \P( Y_{n } \geq \delta) \leq \min_{t\in [0,t_0]} e^{-t \delta} M_{Y_{n} } (t)
.\] 

To find the minimum, we investigate what happens at the extreme points of the interval, namely when
$ t = 0$ and when $ t = \infty$.

\[
    f(0) = e^{-(0)\delta} M_{Y_{n} } (0) = 1
\] 
and
\begin{align*}
    \lim_{t \to \infty} f(t) &= \lim_{t \to \infty} 1/2^{n} e^{-t(\delta + 1/2)} (1+ e^{t/n} )^{n} \\
			     &\geq \lim_{t \to \infty} 1/2^{n} e^{-t(\delta +1/2)} (1 + n (e^{t/n})^{n} )\\
			     &= \lim_{t \to \infty} n/2^{n} e^{-t(\delta - 1/2}  \\
\end{align*}
so if $ \delta \in (0, 1/2)$ then we see that $ e^{-t(\delta - 1/2}$ grows to infinity and so the
limit of the function $ \lim_{t \to \infty} f(t) = \infty$. 
Now let's compare these two extreme values with the value of the function $ f(t)$ at the root of the
derivative $ f'(t)$. Using the result from part (2) of this exercise, we find that the root 
is $ T = n \ln ((1/2 + \delta)/(1/2 - \delta)$ and inserting into $ f(t)$ we get
that $ f(T) = e^{-2n \delta^2} $. Clearly $ f(T) \leq 1, \infty$, and so we see that

\[
    \P(Y_{n} \geq \delta) \leq \min_{t \in [0, \infty)} e^{-t\delta} M_{Y_{n} } (t) = e^{-2n
    \delta^2} 
.\] 

